% !PRINTED 301
\weberpage{301}

Es läßt sich aber $\gamma$ weiter so bestimmen, daß

\weq{12}{A^{\gamma-\lambda}CA^{-\gamma}=C,}

man erhält dafür die Bedingung

\[ \gamma(g-g^{-1})\equiv\lambda g\quad(\mathrm{mod}\ p), \]

die, da $p>3$ ist, immer befriedigt werden kann; danach ist aber wegen (11)

\[ C=A^{\gamma}SA^{-\gamma}. \]

Es kommt also $C$ in dem mit $\mathfrak{K}$ konjugierten Teiler $A^{\gamma}\mathfrak{K} A^{-\gamma}$ vor, und wir können also, indem diese Gruppe an Stelle von $\mathfrak{K}$ gesetzt wird, unbeschadet der Allgemeinheit annehmen, $\mathfrak{K}$ enthalte selbst die Substitution $C$.

Es kann nun in $\mathfrak{K}$ die Substitution

\[ B=\begin{pmatrix}0, & 1\\ -1, & 0\end{pmatrix} \]

entweder vorkommen oder nicht vorkommen. Im ersten Fall enthält $\mathfrak{K}$ auch alle aus $B$ und $C$ zusammengesetzten Substitutionen, die sämtlich von einer der beiden Formen sind:

\weq{13}{\begin{pmatrix}g^{r}, & 0\\ 0, & g^{-r}\end{pmatrix},\qquad r=0,1,\dots\frac{p-3}{2},}

\weq{14}{\begin{pmatrix}0, & g^{-r}\\ -g^{r}, & 0\end{pmatrix},}

und deren Anzahl $p-1$ beträgt. Im ersten Fall enthält also $\mathfrak{K}$ alle Substitutionen von der Form (13), (14), im zweiten alle Substitutionen der Form (13) und \gesperrt{keine} der Form (14).

Es sei nun

\weq{15}{V=\begin{pmatrix}\alpha, & \beta\\ \gamma, & \delta\end{pmatrix}}

eine in $\mathfrak{K}$ enthaltene Substitution, die weder in der Form (13), noch in der Form (14) enthalten ist, bei der also weder $\alpha$ und $\delta$, noch $\beta$ und $\gamma$ zugleich kongruent mit $0$ sind. Dann sind die sämtlichen Substitutionen der Form

\weq{16}{C^{r}VC^{s}=\begin{pmatrix}\alpha g^{r+s}, & \beta g^{r-s}\\ \gamma g^{-r+s}, & \delta g^{-r-s}\end{pmatrix},}

wenn $r,s$ beide die Reihe der Zahlen durchlaufen:

\[ 0,1,2,\dots\frac{p-3}{2} \]

in $\mathfrak{K}$ enthalten und alle \gesperrt{voneinander verschieden}. Denn sind zwei unter den Substitutionen (16) einander gleich, so muß
%CONT>>
